\begin{textsample}{NEG Dim 3 – human – Score 1.00 – sp/sp\_ef\_ciencias\_humanas\_his\_1\_p9.txt}  \label{ex:f3_neg_016}
É também com Bloch e com os Annales que se amplia, já em 1929, a noção de fontes históricas, que passa a ser todo o vestígio deixado pelo homem, desde documentos oficiais a relatos \textbf{orais}.
Atualmente pode¬mos classificar as fontes históricas ou os documentos históricos em : escritas ( cartas, jornais, inventários etc. ) ; \textbf{orais} ( entrevistas, relatos etc. ) ; imateriais ( como festas, saberes populares e tradição \textbf{oral} ) e materiais ( tais como : objetos, construções e indumentária ) ; audiovisuais ( inclusive ficcionais ) ; cartográficas e iconográficas ( pinturas, de¬senhos, fotografias, entre outros ).
As fontes ou documentos históricos são como pistas de um crime para um detetive ou sintomas de uma doença para um médico como nas analogias apresentadas por Carlo Ginzburg sobre o paradigma indiciário, pois para ele : “ O historiador é, por definição, um investigador para quem as experiências, no sentido rigoroso do termo, estão vedadas ” ( GINZBURG, 1989. p.180 ).

% matched lemmas: oral
\end{textsample}
